\chapter{历史高空间案例：统一降为监测，等待重新取证}

本章覆盖的是6只历史高空间样本加3只条件/周期观察样本，和图表1的4只原主标的并不重叠。它们被统一降级并非因为价格不能上涨，而是因为当前价基础的利润分母、现金解释或可审计目标锚至少缺少一项。

\begin{exhibitbox}[图表5｜历史高空间与条件样本的刷新结论]
\small
\begin{tabularx}{\textwidth}{L{1.1cm}L{1.55cm}L{2.1cm}X}
\toprule
\textbf{代码} & \textbf{标的} & \textbf{H1/当前证据} & \textbf{为什么不能恢复旧空间} \\
\midrule
002432 & 九安医疗 & H1预告CNY28--34亿元，科技投资估值变动驱动 [FIN-RAW] & 非经常性投资重估和退出不确定性；无当前公开目标。 \\
000623 & 吉林敖东 & H1预告CNY19.22--21.79亿元，投资收益驱动 [FIN-RAW] & 需要分离医药经营、投资收益和现金。 \\
600739 & 辽宁成大 & H1官方查询缺口 [PROBE-H1]，Q1 OCF为负 & 没有当前H1分母或公开目标锚。 \\
000685 & 中山公用 & H1预告CNY13.31--15.47亿元，投资收益重要 [FIN-RAW] & 当前研报只有预测，无显式目标。 \\
600150 & 中国船舶 & H1官方查询缺口 [PROBE-H1]，Q1 OCF为正 & 仍缺H1交付、回款与当前点目标。 \\
301308 & 江波龙 & 实际H1归母CNY105.77亿元，OCF为负 [F-301308-H1] & 周期峰值、营运资本和EPS分歧尚未归一。 \\
002812 & 恩捷股份 & CNY103外部锚相对当前约84\% [B-002812] & 2027E分母，仍仅为零权重敏感性。 \\
002240 / 300390 & 盛新锂能 / 天华新能 & H1预告改善但Q1 OCF为负 [FIN-RAW] & 周期预告和无显式目标不足以发布范围。 \\
\bottomrule
\end{tabularx}
\normalsize
\sourcenote{FIN-RAW为CNINFO官方财务包；PROBE-H1为600739/600150的官方检索缺口记录；F-301308-H1为2026-08-11公告1225467150第11页；B-002812为东吴证券2026-07-10原始PDF。没有把“未找到”解释为负面事实。}
\end{exhibitbox}

图表5的决策含义是统一的：所有9只只能监测下一次可证实的信息，而不是以历史空间或回撤幅度重新进入正式估值。6只历史高空间样本先满足实际性、现金与锚门槛；3只条件/周期样本还需证明周期恢复可转化为现金。

\section{重新进入估值工作的统一条件}

历史样本要重新获得当前价基础的估值资格，至少需要三件事同时出现：正式中报或季报替代预告/缺口；收入、归母、扣非、OCF和利润质量在同一期间可解释；并取得当前原始目标或可复建的资产/现金流估值锚。任一项缺失都只允许监测，不允许把旧模型空间作为结论。
